\scanpage{300}
\mspage{2\deleted{\unsure{3}}4}
% archive sheet number: 299

$2^\circ$) Est-il encore vrai que si $A$ est dans $\mathrm{Cat}$, $u:F\hookrightarrow G$ un \deleted{flèche}\inserted{monomorphisme} de $\widehat A$, alors $u\in\mathrm{Cof}_{\unsure{W|w}_A}$, et si $u\in W$, $u\in W_A^{\unsure{\mathrm{univ}}}$, donc a-t-on
\[
\mathrm{Mono}_{\widehat A}\subset\mathrm{Cof}_{W_A},\quad
\mathrm{Mono}_{\widehat A}\cap W_A\subset(W_A)^{\mathrm{univ}}\ ?
\]
% slanted
\margin{NB $\mathrm{Loc}(6'a)$ suffit (et $\mathrm{Loc}(6a)$)}

Mais \unsure{$1^\circ$)|ceci} est prouvé \unsure{ds|dans} \uline{XII} prop 20 p 262 -- je viens de vérifier que la démonstr. est valable \deleted{dans} les hypothèses actuelles.$^{\ast}$ \unsure{pourvu|pour} que je vérifie que le cor. 3 p.~149 \uline{XII} est valable, i.e.
% slanted
\margin{Mais elle utilise de plus $\mathrm{Loc}(3\,\mathrm{bis})$. Mais avec \unsure{cette|la} démonstration \uline{plus simple} que je viens de \unsure{rajouter|regarder} (p.~264), on peut s'en passer (et on obtient une vision b. plus complète de la situation)}

$1^\circ$ bis) Si
\diagram{\begin{tikzcd}[ampersand replacement=\&,row sep=small,column sep=small]
X_0\arrow[r,no head]\arrow[d] \& X_1\arrow[d]\\
X_2\arrow[r,no head]\&X
\end{tikzcd}}
est un carré $\mathrm{MV}$ dans $\mathrm{Cat}$, alors
\diagram{\begin{tikzcd}[ampersand replacement=\&,row sep=tiny,column sep=tiny]
\&\&X_1\&\&\\
\displaystyle\int\&X_0\arrow[ur]\arrow[dr]\&\&{}\arrow[r]\&X\\
\&\&X_2\&\&
\end{tikzcd}}
est dans $W$. (donc, moyennant $1^\circ$), le carré est $W$-cocartésien).

Pour le vérifier, je suis amené à vérifier que le th. 4 de loc. cit. (p.~\unsure{146|145}) est encore valable. / je viens de le faire, et ai annoté en conséquence le th. 4 et ses corollaires (certains résultats \inserted{(cor. 4. et cor 5)} utilisent $W(7\,\unsure{bis})$ dans \deleted{toute} sa force, i.e. $\mathrm{Loc}(6')$ i.e. $\mathrm{Loc}(6,a\ \text{\uline{et}}\ b)$. Mais \deleted{\unsure{ces résultats}} cor. 3 du th. 4 est OK.

Je viens aussi de vérifier la validité de $2^\circ$), en vérifiant la dém. du cor. du th 5 (\uline{XII} p.~168) et en l'annotant en conséquence.
